%% USPSC-Cap6-Conclusao.tex
% ----------------------------------------------------------
% Conclusão
% ----------------------------------------------------------
\chapter{Conclusão}
\label{cap:conclusao}

\section{Síntese do Trabalho}
\label{sec:conc_sintese}

% >>> [PREENCHER] Um a dois parágrafos retomando: o problema enunciado na
% Seção 1.2, a abordagem adotada e o percurso realizado. Sem introduzir
% informação nova.

\section{Retomada dos Objetivos}
\label{sec:conc_objetivos}

% >>> [PREENCHER] Percorrer, um a um, os seis objetivos específicos declarados
% na Seção 1.3, indicando onde cada um foi cumprido e com que resultado.
% Se algum objetivo tiver sido cumprido apenas parcialmente, declarar isso
% explicitamente --- a banca valoriza mais a honestidade do que a alegação de
% completude.

\begin{alineas}
  \item \textbf{Caracterizar arquiteturalmente o ecossistema.} \textit{[preencher
  --- remeter ao Quadro do Capítulo~\ref{cap:desenvolvimento}]}

  \item \textbf{Realizar avaliação heurística.} \textit{[preencher]}

  \item \textbf{Definir o perfil de usuário.} \textit{[preencher]}

  \item \textbf{Conduzir avaliação inicial com usuários em campo.}
  \textit{[preencher]}

  \item \textbf{Mensurar a aceitação tecnológica.} \textit{[preencher]}

  \item \textbf{Triangular os resultados.} \textit{[preencher]}
\end{alineas}

\section{Contribuições}
\label{sec:conc_contribuicoes}

% >>> [PREENCHER] Separar em três blocos:
%
% (a) CONTRIBUIÇÃO METODOLÓGICA. A classificação de atritos por origem e o
%     protocolo Think-Aloud híbrido segmentado, propostos para avaliar
%     aplicações com componente de desafio deliberado em contexto de campo.
%     Este é o argumento mais forte do trabalho; desenvolvê-lo com cuidado.
%
% (b) CONTRIBUIÇÃO TÉCNICA. A matriz de rastreabilidade que vincula atrito
%     percebido a decisão de engenharia, e o conjunto priorizado de correções
%     efetivamente implementadas no sistema.
%
% (c) CONTRIBUIÇÃO PARA O PROJETO. O que fica para a equipe do Trilha das
%     Árvores e para a ESALQ: sistema corrigido, instrumentos reaplicáveis e
%     procedimento documentado para rodadas futuras.

\section{Limitações}
\label{sec:conc_limitacoes}

% >>> [PREENCHER] Declarar as limitações sem eufemismo. Entre elas:
% dimensão da amostra e ausência de pretensão inferencial; avaliação restrita
% ao cliente móvel, excluindo o painel administrativo; um único percurso e um
% único dispositivo; papel duplo do pesquisador como desenvolvedor e
% moderador; ausência de avaliação de acessibilidade e de qualidades hedônicas.

\section{Considerações Sobre o Curso de Graduação}
\label{sec:conc_graduacao}

% >>> [PREENCHER] Seção de caráter pessoal, prevista pelo formato do TCC.
% Sugestão de eixos: quais disciplinas do curso sustentaram o trabalho; que
% lacunas de formação foram percebidas ao longo dele --- em especial na
% articulação entre engenharia de software e métodos de pesquisa com pessoas;
% e o que a experiência de conduzir um estudo com usuários acrescentou à
% formação como engenheiro.

\section{Trabalhos Futuros}
\label{sec:conc_futuros}

% >>> [PREENCHER] Derivar cada item de algo efetivamente identificado no
% trabalho, e não de uma lista genérica. Candidatos já levantados ao longo
% do texto:
%
%   - Conclusão da integração do serviço de armazenamento de objetos e
%     remoção dos artefatos binários do banco relacional.
%   - Adoção de TLS e revisão do modelo de emissão de tokens (expiração e
%     gestão de chave).
%   - Operação com base cartográfica disponível sem conectividade.
%   - Integração da tela de perfil a dados reais.
%   - Atualização da especificação da API.
%   - Aplicação do instrumento a amostra ampliada, junto a turmas da ESALQ,
%     com verificação de consistência interna dos construtos.
%   - Investigação das qualidades hedônicas da experiência por instrumentos
%     como o UEQ ou o AttrakDiff.
%   - Avaliação de acessibilidade segundo a WCAG.
%   - Avaliação do painel administrativo com o perfil de gestor.
